\section{Thiết kế Monitoring và cảnh báo}
\label{sec:monitoring-design}

Giám sát được thiết kế theo hai kênh bổ trợ nhau. Kênh \textbf{telemetry ứng
dụng} đi qua MQTT, mang ngữ nghĩa của nền tảng (runtime, FPS, camera, chất lượng
đầu ra model) và phục vụ Dashboard, Assistant. Kênh \textbf{metric hạ tầng} dùng
hệ sinh thái Prometheus \cite{prometheus} để thu metric node, container và GPU một cách chuẩn hóa,
lưu dài hạn và cảnh báo theo luật.

\subsection{Đặc tả telemetry}

Đặc tả telemetry chia chỉ số thành bảy nhóm và ba loại thông điệp
(Bảng~\ref{tab:telemetry-groups}). Tất cả trường đều tùy chọn để thiết bị chỉ báo
cáo những gì đo được; tên trường thống nhất giữa Edge Agent, kho lưu trữ, bộ mô
phỏng dữ liệu và pipeline học máy.

\begin{table}[H]
\centering
\caption{Nhóm đặc trưng telemetry của Edge Node}
\label{tab:telemetry-groups}
\begin{tabularx}{\textwidth}{L{3cm}L{2.2cm}Y}
\toprule
\textbf{Nhóm} & \textbf{Thông điệp} & \textbf{Ví dụ trường}\\
\midrule
Tài nguyên node & system & \texttt{cpu\_usage\_pct}, \texttt{cpu\_load\_1m/5m}, \texttt{memory\_usage\_pct}, \texttt{swap\_usage\_pct}\\
Nhiệt và điện & system, gpu & \texttt{temperature\_c}, \texttt{thermal\_throttling}, \texttt{power\_w}, \texttt{gpu\_temperature\_c}\\
Lưu trữ & system & \texttt{disk\_usage\_pct}, \texttt{disk\_free\_gb}, \texttt{inode\_usage\_pct}\\
Mạng & system & \texttt{network\_rx/tx\_mb\_s}, \texttt{packet\_loss\_pct}, \texttt{mqtt\_publish\_latency\_ms}\\
GPU & gpu & \texttt{gpu\_usage\_pct}, \texttt{gpu\_memory\_used\_mb}, \texttt{gpu\_frequency\_mhz}, \texttt{gpu\_power\_w}\\
Runtime và suy luận & runtime & \texttt{inference\_fps}, \texttt{effective\_fps\_ratio}, \texttt{inference\_p95\_ms}, \texttt{inference\_error\_rate}, \texttt{runtime\_queue\_size}\\
Camera và đầu ra model & runtime & \texttt{camera\_fps}, \texttt{camera\_frame\_drop\_pct}, \texttt{camera\_freeze\_duration\_sec}, \texttt{detection\_count}, \texttt{confidence\_mean}, \texttt{class\_entropy}\\
\bottomrule
\end{tabularx}
\end{table}

Telemetry \texttt{system} được thu bằng \texttt{psutil}; telemetry \texttt{gpu} được
thu bằng \texttt{nvidia-smi} trên máy có GPU rời (collector trả về rỗng khi công
cụ không có, ví dụ trên Jetson); telemetry \texttt{runtime} do AI Runtime tự đo.

Nhóm \emph{đầu ra model} có ý nghĩa đặc biệt: \texttt{confidence\_mean},
\texttt{detection\_count} và \texttt{class\_entropy} cho phép nhận biết suy giảm
chất lượng suy luận (camera mờ, bị che, đổi góc, ánh sáng thay đổi, dữ liệu trôi)
ngay cả khi CPU, GPU, FPS và latency vẫn bình thường --- điều mà giám sát hạ tầng
thuần túy không phát hiện được.

\subsection{Metric hạ tầng với Prometheus}

Trên mỗi Edge Node, \texttt{node-exporter} (metric hệ điều hành),
\texttt{cAdvisor} (metric container) và \texttt{dcgm-exporter} (GPU NVIDIA, tùy
chọn) chỉ lắng nghe trên \texttt{127.0.0.1}. Một Prometheus chạy ở
\emph{agent mode} thu các exporter mỗi 15~s, gắn nhãn \texttt{node\_id},
\texttt{cluster\_id} và đẩy về Prometheus trên Cloud bằng \texttt{remote\_write}
có xác thực. Mô hình đẩy (push) hoạt động được khi Edge nằm sau NAT; WAL của
agent đệm mẫu khi mất mạng và gửi bù khi kết nối lại.

Trên Cloud, Prometheus lưu metric (mặc định 30 ngày) và đánh giá luật cảnh báo;
Alertmanager định tuyến cảnh báo; Grafana hiển thị và cho phép nhúng panel vào
Dashboard. Bộ luật cảnh báo Edge được liệt kê trong Bảng~\ref{tab:alert-rules}.

\begin{table}[H]
\centering
\caption{Luật cảnh báo metric hạ tầng Edge}
\label{tab:alert-rules}
\begin{tabularx}{\textwidth}{L{4.2cm}Y C{1.6cm}C{1.8cm}}
\toprule
\textbf{Cảnh báo} & \textbf{Điều kiện} & \textbf{Kéo dài} & \textbf{Mức}\\
\midrule
EdgeNodeDown & Node từng báo cáo trong 1 ngày nhưng không đẩy metric trong 5 phút & 2 phút & critical\\
EdgeDiskAlmostFull & Phân vùng gốc sử dụng trên 90\% & 10 phút & warning\\
EdgeHighMemory & Bộ nhớ sử dụng trên 90\% & 10 phút & warning\\
EdgeContainerRestarting & Container khởi động lại hơn 2 lần trong 15 phút & -- & warning\\
EdgeGpuHot & Nhiệt độ GPU trên 85\,°C & 5 phút & warning\\
\bottomrule
\end{tabularx}
\end{table}

\subsection{Thiết kế mô-đun cảnh báo sớm}

Trên nền dữ liệu telemetry đã chuẩn hóa, đồ án thiết kế mô-đun cảnh báo sớm
suy giảm (early warning) theo luồng ở Hình~\ref{fig:monitoring-incident-flow}.

\begin{figure}[H]
\centering
\resizebox{\textwidth}{!}{%
\begin{tikzpicture}[node distance=0.6cm and 0.6cm]
  \node[flowbox] (t) {Telemetry\\5 s/mẫu};
  \node[flowbox, right=of t] (s) {Lưu trữ\\chuỗi thời gian};
  \node[flowbox, right=of s, fill=orange!8] (w) {Cửa sổ trượt\\5 phút (60 mẫu)};
  \node[flowbox, right=of w] (f) {Đặc trưng thời gian};
  \node[flowbox, right=of f, fill=green!6] (m) {Isolation Forest\\XGBoost};
  \node[flowbox, right=of m, fill=red!6] (a) {Risk score\\Cảnh báo};
  \draw[flowarrow] (t) -- (s);
  \draw[flowarrow] (s) -- (w);
  \draw[flowarrow] (w) -- (f);
  \draw[flowarrow] (f) -- (m);
  \draw[flowarrow] (m) -- (a);
\end{tikzpicture}}
\caption{Luồng xử lý của mô-đun cảnh báo sớm suy giảm}
\label{fig:monitoring-incident-flow}
\end{figure}

\textbf{Đặc trưng thời gian.} Với chu kỳ lấy mẫu 5~s và cửa sổ 5 phút, mỗi chỉ
số cốt lõi được biến đổi thành \texttt{current}, \texttt{mean}, \texttt{min},
\texttt{max}, \texttt{std}, \texttt{slope} (hệ số hồi quy tuyến tính) và
\texttt{delta} (chênh lệch so với đầu cửa sổ). Bộ chỉ số khởi đầu gồm khoảng 20
trường: CPU, bộ nhớ và tốc độ tăng bộ nhớ, GPU, nhiệt độ và tốc độ tăng nhiệt, đĩa,
RTT, mất gói, FPS và \texttt{effective\_fps\_ratio}, p95 latency, tỷ lệ lỗi, kích
thước hàng đợi, số lần restart, FPS camera, tuổi khung, tỷ lệ drop, số detection và
confidence trung bình. Một số đặc trưng dẫn xuất như
\texttt{latency\_ratio} = p95 hiện tại / p95 baseline và tốc độ tăng hàng đợi giúp
mô hình học quan hệ đa biến thay vì ngưỡng cứng.

\textbf{Nhãn trạng thái.} Mỗi thời điểm được gắn một trong bốn trạng thái vận
hành: NORMAL, WARNING, DEGRADED, CRITICAL. Ví dụ DEGRADED khi FPS giảm rõ, latency
và hàng đợi tăng; CRITICAL khi runtime FAILED, camera mất kết nối hoặc node OFFLINE.

\textbf{Hai bài toán tách biệt.} Isolation Forest trả lời ``hành vi hiện tại có bất
thường không'' (không cần nhãn). XGBoost trả lời ``trong 10 phút tới có suy giảm
không'': nhãn \texttt{future\_degradation}=1 nếu node chuyển sang DEGRADED hoặc
CRITICAL trong 10 phút sau cửa sổ quan sát. Để tránh rò rỉ nhãn, các trường chỉ
xuất hiện sau sự cố (ví dụ \texttt{runtime\_status}=FAILED) không được dùng làm
đầu vào dự báo. Dữ liệu thô luôn được giữ lại để có thể thay đổi cách trích đặc
trưng mà không cần thu thập lại.

Ví dụ phân biệt: khi GPU tăng, nhiệt độ tăng, latency và hàng đợi tăng trong khi
FPS và confidence giảm, mô hình cảnh báo nguy cơ suy giảm dù từng chỉ số chưa
vượt ngưỡng. Ngược lại, khi GPU tăng do cảnh có nhiều đối tượng nhưng latency,
FPS, hàng đợi và confidence ổn định, trạng thái vẫn là NORMAL, tránh cảnh báo sai.
Trong phạm vi đồ án, hợp đồng telemetry, thu thập và lưu trữ đã được hiện thực;
việc huấn luyện và tích hợp mô hình cảnh báo sớm vào luồng cảnh báo được trình
bày ở mức thiết kế và là hướng phát triển tiếp theo.
